\documentclass{article}
\usepackage{amsfonts,amsthm,amsmath}
\usepackage[mathscr]{euscript}
\usepackage{enumitem}

\setcounter{section}{0}
\renewcommand{\thesection}{\S\arabic{section}}
\renewcommand{\thesubsection}{\arabic{section}}

\newtheorem{thm}{Theorem}
\newenvironment{theorem}[1]{
	\renewcommand{\thethm}{#1}
	\begin{thm}
		}{\end{thm}}

\newtheorem{lma}{Lemma}
\newenvironment{lemma}[1]{
	\renewcommand{\thelma}{#1}
	\begin{lma}
		}{\end{lma}}

\theoremstyle{definition}
\newtheorem{ex}{Exercise}
\newenvironment{exercise}[1]{
	\renewcommand{\theex}{#1}
	\begin{ex}
		}{\end{ex}}

\title{Exercises 8.B}
\author{Connor Mooneyhan}
\date{February 25, 2024}

\begin{document}

\maketitle

\begin{ex}
	Suppose $V$ is a complex vector space, $N \in \mathcal{L}(V)$, and 0 is the only eigenvalue of $N$. Prove that $N$ is nilpotent.
\end{ex}
\begin{proof}
	By 8.21, \begin{align*}
		V & = G(0, N)                                   \\
		  & = \mathrm{null}\,(N - 0I)^{\mathrm{dim}\,V} \\
		  & = \mathrm{null}\,N^{\mathrm{dim}\,V},
	\end{align*}
	so given $v \in V$, we have $N^{\mathrm{dim}\,V}v = 0$, thus $N$ is nilpotent.
\end{proof}

\begin{ex}
	Give an example of an operator $T$ on a finite-dimensional real vector space such that 0 is the only eigenvalue of $T$ but $T$ is not nilpotent.
\end{ex}
\begin{proof}
	Let $T \in \mathbb{R}^3$ be defined by $T(z_1, z_2, z_3) = (-z_2, z_1, 0)$. Then the only eigenvalue of $T$ is $0$, but $T^3(z_1, z_2, z_3) = (z_2, -z_1, 0)$, so $T$ is not nilpotent.
\end{proof}

\begin{ex}
	Suppose $T \in \mathcal{L}(V)$. Suppose $S \in \mathcal{L}(V)$ is invertible. Prove that $T$ and $S^{-1}TS$ have the same eigenvalues with the same multiplicities.
\end{ex}
\begin{proof}
	We seek to show that $S|_{G(\lambda, S^{-1}TS)}$ is in fact an isomorphism from $G(\lambda, S^{-1}TS)$ to $G(\lambda, T)$ for $\lambda \in \mathbf{F}$, which shows that when the former is nontrivial (so that $\lambda$ is an eigenvalue of $S^{-1}TS$), so is the latter, and vice versa.
	We already know that $S|_{G(\lambda, S^{-1}TS)}$ is invertible since $S$ is invertible, so it is an isomorphism onto its range.
	Thus we need only show that $G(\lambda, T) = \mathrm{range}\,S|_{G(\lambda, S^{-1}TS)}$.

	We know that $(S^{-1}TS - \lambda I)^{\mathrm{dim}\,V} = S^{-1}(T - \lambda I)^{\mathrm{dim}\,V}S$ since $p(S^{-1}TS) = S^{-1}p(T)S$ for any $p, T,$ and invertible $S$.
	For a given eigenvalue $\lambda$ of $S^{-1}TS$, let $v \in G(\lambda, S^{-1}TS)$.
	Then \begin{align*}
		0 & = (S^{-1}TS - \lambda I)^{\mathrm{dim}\,V}v \\
		  & = S^{-1}(T - \lambda I)^{\mathrm{dim}\,V}Sv \\
		  & = (T - \lambda I)^{\mathrm{dim}\,V}Sv,
	\end{align*}
	where the final equality comes from applying $S$ to both sides on the left.
	Thus $Sv \in G(\lambda, T)$.

	Likewise, given $u \in G(\lambda, T)$, we can express $u$ uniquely as $u = Sv$ for some $v \in V$ since $S$ is surjective, so \begin{align*}
		0 & = (T - \lambda I)^{\mathrm{dim}\,V}u         \\
		  & = (T - \lambda I)^{\mathrm{dim}\,V}Sv        \\
		  & = S^{-1}(T - \lambda I)^{\mathrm{dim}\,V}Sv  \\
		  & = (S^{-1}TS - \lambda I)^{\mathrm{dim}\,V}v,
	\end{align*}
	where the third equality comes from applying $S^{-1}$ to both sides on the left.
	Thus $v \in G(\lambda, S^{-1}TS)$.

	This means $S$ is an isomorphism from $G(\lambda, S^{-1}TS)$ to $G(\lambda, T)$, so \begin{equation*}
		\mathrm{dim}\,G(\lambda, S^{-1}TS) = \mathrm{dim}\,G(\lambda, T).
	\end{equation*}
\end{proof}

\begin{ex}
	Suppose $V$ is an $n $-dimensional complex vector space and $T$ is an operator on $V$ such that $\mathrm{null}\,T^{n - 2} \neq \mathrm{null}\,T^{n - 1}$. Prove that $T$ has at most two distinct eigenvalues.
\end{ex}
\begin{proof}

\end{proof}

\end{document}
